\documentclass[10pt,a4paper]{amsart}
\usepackage[margin=19mm]{geometry}
\usepackage{amsmath,amssymb,mathtools}
\usepackage[scr=rsfs,cal=euler]{mathalfa}
\usepackage{xfrac}
\usepackage[unicode,hidelinks,bookmarks=false]{hyperref}
\newcommand{\ie}{i.e.\ }
\newcommand{\cf}{cf.\ }
\newcommand{\resp}{respectively\ }
\newcommand{\Subsection}{\subsection}
\providecommand{\nobreakdash}{\nobreak\hbox{-}}
\setlength{\emergencystretch}{2em}
\multlinegap0pt
\usepackage{amssymb}
\usepackage{booktabs}
\usepackage{caption}
\newtheorem{prop}{Proposition}
\title{Complex quaternionic manifolds and c-projective structures}
\author{Aleksandra Bor\'owka}
\date{Source: arXiv:2412.10301v2 (CC BY 4.0)}
\begin{document}
\maketitle

\noindent\emph{Source and licence.} Complex quaternionic manifolds and c-projective structures. Version arXiv:2412.10301v2 (CC BY 4.0), distributed by arXiv under the Creative Commons Attribution 4.0 International licence, http://creativecommons.org/licenses/by/4.0/, as recorded in the arXiv licence metadata for this exact version and shown on the abstract page https://arxiv.org/abs/2412.10301v2. Full licence notice: \texttt{licenses/arxiv-2412.10301-CC-BY-4.0.txt}.\par\smallskip

\noindent\textit{Editorial scope.} Counted unit: the article's prop con (source0.tex lines 250--252). The article's complete original proof of it is delivered with it (source0.tex lines 253--255, one \texttt{proof} environment of 3 lines). No mathematics is written, repaired or re-derived: the statement and the proof are the article's own lines, in the article's own notation, and no retained line is substituted or re-flowed.  Retained uncounted as the source-local context the proof needs: lines 240--242.  Nothing was omitted from any retained span, so the package carries no omission record.  No retained line contains a citation, so the wrapper carries no bibliography and no bibliography entry.  No retained line contains a figure environment, an image-inclusion command or drawing code, so no image is delivered and the package declares no asset dependency.  Editorial formatting only, in addition: an AMS \texttt{amsart} preamble that restates the article's own declarations and macros used by the retained lines, namely the environments \texttt{prop} on separate counters, together with the standard packages those lines need.  The retained environments print in this document as Proposition 1 in order of appearance, whereas the article prints the same items with the numbering of its own full document.  The complete native source of the article as distributed in the arXiv e-print for this exact version is bundled unchanged as \texttt{sources/170081/source0.tex}.
\begin{prop}\label{main}
 The hypersurfaces $D^+\cap U^+\subset Z$ and $D^-\cap U^-\subset Z$ define up to a sign an integrable $S^1$-invariant complex structure on the quaternionic manifold $M$ obtained from the twistor space $Z$.
\end{prop}

\begin{prop}\label{con}
    The $GL(n,\mathbb{H})U(1)$ connection corresponding to the complex structure obtained in  Proposition \ref{main} is invariant with respect to the $S^1$-action.
\end{prop}

\begin{proof}
  Let $\mathcal{D}^J$ be the $GL(n,\mathbb{H})U(1)$ connection with $\mathcal{D}^J J_D=0$. We know that for a generator $X$ of the action the push forward by the flow of $\mathcal{D}^J$ along $X$ by any fixed time  is some quaternionic connection $\mathcal{D}'$ (since $X$ preserves the quaternionic structure). Moreover, $\mathcal{D}'$ preserves the underlying push forward of $J$.  But as $L_XJ=0$, we get that  $\mathcal{D}'J=0$. As  $\mathcal{D}^J$ is the unique connection with  $\mathcal{D}^J J_D=0$, we get that  $\mathcal{D}^J =\mathcal{D}'$. 
\end{proof}

\end{document}
